% Shared-T5 Multi-DocID method section.
% This file is intended to be included from the main manuscript.

\section{Method}
\label{sec:method}

\subsection{Problem Formulation}
\label{sec:problem}

Let $\mathcal{D}=\{d_1,\ldots,d_N\}$ be a fixed corpus and let $q$ denote a query with a set of relevant documents $\mathcal{D}^{+}(q)\subseteq\mathcal{D}$. Generative retrieval replaces explicit query--document similarity search with autoregressive identifier generation~\citep{tay2022dsi}. A model maps $q$ to one or more discrete routes, which are subsequently resolved to documents in the corpus.

A single identifier can constitute a brittle retrieval path: documents that share an identifier are indistinguishable under that path, while an incorrectly generated early token can remove the relevant document from the beam. We therefore assign each document $V=3$ identifiers and retrieve through all three routes using one shared sequence-to-sequence model. We refer to this construction as \emph{Multi-DocID}. The method consists of four stages:
\begin{enumerate}[leftmargin=*,nosep]
    \item construct a nine-level residual-quantized code for each document;
    \item partition the code into three residual-code routes;
    \item train one shared T5 to generate every route under explicit view conditioning; and
    \item run three constrained decodes and fuse their evidence at the document level.
\end{enumerate}
Although we use \emph{view} as a concise index for the three routes, these are disjoint slices of one residual-quantization chain rather than independent modalities.

\subsection{Residual-Quantized Multi-DocIDs}
\label{sec:rq-multidocid}

\subsubsection{Nine-level residual quantization}

We first encode every document into a fixed dense representation
\begin{equation}
    \mathbf{x}_d=f_{\mathrm{doc}}(d)\in\mathbb{R}^{D}.
\end{equation}
To obtain a discrete identifier, we train an $L$-level residual quantizer with $L=9$. For each level $\ell\in\{0,\ldots,L-1\}$, let
\begin{equation}
    \mathcal{E}_{\ell}=\{\mathbf{e}_{\ell,k}\}_{k=0}^{K-1}
\end{equation}
be a codebook with $K=256$. Starting with $\mathbf{r}^{(0)}_d=\mathbf{x}_d$, greedy residual quantization selects
\begin{align}
    z_{d,\ell}
    &=\arg\min_{k}\left\|\mathbf{r}^{(\ell)}_d-\mathbf{e}_{\ell,k}\right\|_2^2, \\
    \mathbf{r}^{(\ell+1)}_d
    &=\mathbf{r}^{(\ell)}_d-\mathbf{e}_{\ell,z_{d,\ell}}.
\end{align}
The reconstructed embedding is
\begin{equation}
    \widehat{\mathbf{x}}_d
    =\sum_{\ell=0}^{L-1}\mathbf{e}_{\ell,z_{d,\ell}}.
\end{equation}
Following residual-quantized Semantic IDs~\citep{lee2022rq,tiger2023}, the implemented quantizer minimizes
\begin{align}
    \mathcal{L}_{\mathrm{RQ}}
    ={}&\left\|\mathbf{x}_d-\widehat{\mathbf{x}}_d\right\|_2^2 \\
    &+\frac{\lambda_{\mathrm{cb}}}{L}
      \sum_{\ell=0}^{L-1}
      \left\|\operatorname{sg}[\mathbf{r}^{(\ell)}_d]
      -\mathbf{e}_{\ell,z_{d,\ell}}\right\|_2^2 \\
    &+\frac{\lambda_{\mathrm{com}}}{L}
      \sum_{\ell=0}^{L-1}
      \left\|\mathbf{r}^{(\ell)}_d
      -\operatorname{sg}[\mathbf{e}_{\ell,z_{d,\ell}}]\right\|_2^2,
\end{align}
where $\operatorname{sg}[\cdot]$ denotes stop-gradient; averaging over documents in a minibatch is implicit. We use $\lambda_{\mathrm{cb}}=1.0$ and $\lambda_{\mathrm{com}}=0.25$. We initialize each level by greedy $k$-means on its input residuals and periodically re-seed unused codes. Unlike an RQ-VAE, our quantizer operates directly on fixed document embeddings and does not learn an additional encoder--decoder around them. The implementation retains the conventional commitment term in the logged objective; because the document embeddings and residual inputs are fixed or detached, this term does not contribute gradients to the trainable quantizer parameters.

\subsubsection{Three residual-code routes}

The full RQ output is a nine-code tuple
\begin{equation}
    \mathbf{z}(d)=(z_{d,0},z_{d,1},\ldots,z_{d,8}).
\end{equation}
Rather than decode this tuple as one nine-token identifier, we split it into exactly three contiguous routes:
\begin{align}
    \rho_0(d)&=(z_{d,0},z_{d,1},z_{d,2}),\\
    \rho_1(d)&=(z_{d,3},z_{d,4},z_{d,5}),\\
    \rho_2(d)&=(z_{d,6},z_{d,7},z_{d,8}).
\end{align}
The routes have a global level-specific token namespace:
\begin{align}
    \rho_0(d)&\mapsto \langle r0\_z\rangle\ \langle r1\_z\rangle\ \langle r2\_z\rangle,\\
    \rho_1(d)&\mapsto \langle r3\_z\rangle\ \langle r4\_z\rangle\ \langle r5\_z\rangle,\\
    \rho_2(d)&\mapsto \langle r6\_z\rangle\ \langle r7\_z\rangle\ \langle r8\_z\rangle.
\end{align}
Here the placeholder $z$ is replaced by the selected code index. Thus, codes with the same index at different RQ levels remain distinct tokens. We add all $LK$ route tokens and three view markers $\langle\mathrm{view}_v\rangle$ to the T5 vocabulary and require each to be atomic under tokenization.

The split exposes earlier, middle, and later stages of the residual chain: each later stage quantizes what remains after the preceding codewords have been subtracted. The quantizer does not, however, impose monotonic codeword magnitudes or guarantee a coarse-to-fine semantic hierarchy across these slices. Our hypothesis is operational: their different collision patterns and query-conditioned decoding errors can provide complementary retrieval paths.

\subsubsection{Collision-preserving route index}

A three-token route is not required to uniquely identify a document. For each view $v$, we retain the complete posting list
\begin{equation}
    \mathcal{P}_v(r)=\{d\in\mathcal{D}:\rho_v(d)=r\}.
\end{equation}
This differs from appending an arbitrary uniqueness suffix, as used in some RQ-based generative recommenders~\citep{tiger2023}: a suffix removes collisions but introduces an additional token whose value may not be semantically predictable from a query. Our design preserves the ambiguity rather than dropping colliding owners. During training, documents that collide in a view intentionally share that view's target route. At inference, we allocate that route's evidence across all owners, while alternate routes may provide evidence that differentiates them. Each view-specific trie stores the valid three-token route language; a separate per-view posting index maps each decoded leaf route to $\mathcal{P}_v(r)$.

\subsection{Shared-T5 Training}
\label{sec:shared-training}

\subsubsection{Indexing and retrieval supervision}

Following the indexing/retrieval decomposition of DSI~\citep{tay2022dsi}, we train on both document content and query supervision. Let $\mathcal{X}(d)$ denote the textual sources associated with document $d$. It includes the document text, labeled training queries, and, when used, document-conditioned pseudo queries generated by docT5query~\citep{nogueira2019doct5query}. We prefix document and query sources with the strings \texttt{document:} and \texttt{query:}, respectively.

For every source $x\in\mathcal{X}(d)$, we construct one example per route:
\begin{equation}
    \left(\langle\mathrm{view}_v\rangle\oplus x,\ \rho_v(d)\right),
    \qquad v\in\{0,1,2\},
\end{equation}
where $\oplus$ denotes textual concatenation. For example,
\begin{align*}
    \text{input: }&\langle\mathrm{view}_1\rangle\ \texttt{query: }q,\\
    \text{target: }&\langle r3\_a\rangle\ \langle r4\_b\rangle\ \langle r5\_c\rangle.
\end{align*}
The view marker occurs only in the input; the target contains the corresponding three route tokens followed by EOS.

A single T5 model~\citep{raffel2020t5} with parameters $\theta$ is shared across all sources and views. Its training objective is standard teacher-forced sequence cross-entropy:
\begin{equation}
\label{eq:gr-loss}
    \mathcal{L}_{\mathrm{GR}}
    =-\sum_{d\in\mathcal{D}}
      \sum_{x\in\mathcal{X}(d)}
      \sum_{v=0}^{V-1}
      \sum_{t=1}^{|\rho_v(d)|+1}
      \log p_{\theta}\!\left(
        y_{d,v,t}\mid y_{d,v,<t},
        \langle\mathrm{view}_v\rangle\oplus x
      \right).
\end{equation}
Unlike a three-model ensemble, this expansion does not multiply model parameters: all route mappings are learned in the same encoder--decoder.

\subsubsection{Geometry-aware route-token initialization}

At the beginning of a new run, we initialize the embedding of route token $\langle r\ell\_k\rangle$ with its corresponding normalized RQ centroid. Let $\overline{n}$ be the mean norm of the pretrained T5 vocabulary embeddings. We set
\begin{equation}
    \mathbf{w}_{\ell,k}
    \leftarrow
    \alpha\,\overline{n}\,
    \frac{\mathbf{e}_{\ell,k}}
         {\|\mathbf{e}_{\ell,k}\|_2},
\end{equation}
with scale $\alpha=1$ in our experiments. T5 ties its input and output embeddings; therefore, the route logits are initialized consistently with the quantizer geometry instead of treating all new route tokens as unrelated random labels. The view-marker embeddings retain their standard newly-added-token initialization.

\subsection{Multi-Route Constrained Retrieval}
\label{sec:inference}

\subsubsection{View-conditioned constrained decoding}

At inference time, the same query is decoded once per view:
\begin{equation}
    q_v=\langle\mathrm{view}_v\rangle\oplus\texttt{query: }q.
\end{equation}
For each $q_v$, the shared T5 performs beam search under the corresponding view-specific trie. The trie constrains every prefix to a route observed in the corpus and forces the correct global level sequence. Let the raw beam output be
\begin{equation}
    \widetilde{\mathcal{B}}_v(q)
    =\bigl[(r_{v,b},s_{v,b})\bigr]_{b=1}^{B},
\end{equation}
where $s_{v,b}$ is the length-normalized sequence score. We deduplicate identical route strings within a view and define
\begin{align}
    \mathcal{B}_v(q)
    &=\{r_{v,b}:1\leq b\leq B\},\\
    s_v(r\mid q)
    &=\max_{b:r_{v,b}=r}s_{v,b}.
\end{align}
Thus $\mathcal{B}_v(q)$ is a set of at most $B$ unique routes and $s_v$ retains the best score for each route. Our primary operating point allocates the full raw beam size $B$ to every view, rather than dividing a fixed total beam budget across views.

The document candidate set is the union of the returned posting lists:
\begin{equation}
\label{eq:candidate-union}
    \mathcal{C}(q)
    =\bigcup_{v=0}^{V-1}
      \bigcup_{r\in\mathcal{B}_v(q)}
      \mathcal{P}_v(r).
\end{equation}
A document retrieved through multiple routes is represented once in $\mathcal{C}(q)$ while retaining all of its per-view evidence.

\subsubsection{Beam-normalized route weights}

Raw sequence scores are not directly comparable across views. Our primary fusion rule first normalizes them within each returned view beam:
\begin{equation}
\label{eq:route-weight}
    \widehat{w}_{v,B}(r\mid q)
    =
    \frac{\exp s_v(r\mid q)}
         {\sum_{r'\in\mathcal{B}_v(q)}
          \exp s_v(r'\mid q)}.
\end{equation}
Because $s_v$ is Hugging Face's length-normalized \texttt{sequences\_scores}, Equation~\ref{eq:route-weight} defines a normalized weight over the returned beams, not an exact model posterior over the trie. All valid routes have the same generated length, so the transformation is applied consistently within a view, but the weights remain conditional on the truncated returned set.

We distribute each route's mass uniformly among the documents in its posting list:
\begin{equation}
\label{eq:view-doc-weight}
    \widehat{w}_{v,B}(d\mid q)
    =
    \sum_{\substack{r\in\mathcal{B}_v(q):\\d\in\mathcal{P}_v(r)}}
    \frac{\widehat{w}_{v,B}(r\mid q)}
         {|\mathcal{P}_v(r)|}.
\end{equation}
The uniform division is an explicit allocation assumption for otherwise indistinguishable owners; it prevents a large posting list from assigning the full route weight to every document.

Finally, we average document evidence across views:
\begin{equation}
\label{eq:beam-weight-fusion}
    S_{\mathrm{BW}}(d,q)
    =\frac{1}{V}\sum_{v=0}^{V-1}
      \widehat{w}_{v,B}(d\mid q).
\end{equation}
For numerical stability, the implementation ranks by $\log S_{\mathrm{BW}}$, computed with log-sum-exp; this is ranking-equivalent to Equation~\ref{eq:beam-weight-fusion}. Ties are broken by the best route rank and then the document index. This fusion is parameter-free, operates after route-to-document expansion, and assigns equal weight to each decoding view. We use \emph{beam posterior} only as an implementation shorthand for this beam-normalized weighting rule.

\subsection{Alternative Document-Level Fusion Rules}
\label{sec:fusion-ablations}

We retain three alternative families to separate candidate complementarity from the effect of score aggregation. These are diagnostic alternatives rather than the primary inference rule in Equation~\ref{eq:beam-weight-fusion}.

\paragraph{Majority voting.}
MINDER establishes the broader pattern of aggregating evidence from multiple generated identifiers~\citep{li2023minder}. As a deliberately confidence-free baseline, we rank documents by cross-view support:
\begin{equation}
    S_{\mathrm{vote}}(d,q)
    =\sum_{v=0}^{V-1}
      \mathbb{I}\!\left[
        d\in\bigcup_{r\in\mathcal{B}_v(q)}\mathcal{P}_v(r)
      \right].
\end{equation}
Ties are resolved by the best route rank across views and then the document index. This provides a support-only ranking baseline that discards route-score magnitude and collision-size correction.

\paragraph{Validation-calibrated score fusion.}
For each view, we fit a temperature $T_v>0$ and bias $b_v$ by binary cross-entropy on validation candidate--view scores with document-relevance labels:
\begin{equation}
    g_v(s)=s/T_v+b_v.
\end{equation}
The calibrated LogSumExp score is
\begin{equation}
    S_{\mathrm{LSE}}(d,q)
    =\log\sum_{v\in\mathcal{H}(d,q)}
      \exp g_v(s_v(d\mid q)),
\end{equation}
where $\mathcal{H}(d,q)$ is the set of views that retrieved $d$, and $s_v(d\mid q)$ is the best route score through which view $v$ retrieved it. Sum and max aggregation are corresponding ablations. Calibration parameters are fit only on validation data and then frozen for test evaluation.

\paragraph{Learned fusion diagnostic.}
We additionally train a linear logistic document reranker using 17 features: for each of three views, calibrated score, presence, reciprocal route rank, and $\log(1+\text{collision size})$; plus hit count, best reciprocal rank, and calibrated LSE, max, and mean. For a missing view, its calibrated-score feature is imputed with that view's minimum observed validation score, while its presence, rank, and collision features are zero. The class-balanced model is fit on validation candidates with selected hard negatives. It diagnoses learnable fusion headroom but is not a mathematical upper bound; because it introduces learned parameters and a separate reranking stage, we do not treat it as the main Multi-DocID method.

\subsection{Computational Characteristics}
\label{sec:complexity}

Multi-DocID retains one T5 parameter set and one tokenizer, but performs $V=3$ constrained decodes. With beam size $B$ per view and fixed three-token routes, it returns at most $VB$ route hypotheses before posting-list expansion. Trie traversal and posting expansion are view-specific, after which scoring is document-level. We report the three-decode latency and compute explicitly. To distinguish route diversity from search-width effects, our controlled comparison increases the beam of a single view to approximately match the total number of returned routes; it is a search-budget control rather than a claim of equal wall-clock latency.
